\thispagestyle{empty}
{\Huge Prefácio}

\vspace{1cm}
%%%%%%%%%%%%%%%%%%%%%%%%%%%%%%%%%%%%%%%%%%%%%%%%%%%%%%%%%%%%%%%%%%%%%
Com este texto, inauguro uma coleção de textos que gostaria de ter
iniciado muito antes. Uma coleção chamada {\it Lições Curtas em
Física Universitária}. O objetivo desta coleção é ajudar na formação
de alunos da área de ciências exatas que não possuem o domínio da 
lingua inglesa, e tem dificuldades em ler textos em inglês. A grande
maioria dos textos técnicos são publicados em inglês. O conteúdo
dos textos dessa coleção serão bem vastos, abrangendo desde assuntos
de início de graduação, como por exemplo física básica, termodinâmica,
eletromagnetismo, e até mesmo temas de pesquisa, como Efeito Kondo,
Cosmologia, Relatividade Geral, etc.

Não pretendo fazer mais do mesmo. Como a maioria dos autores fazem,
ou por vício, ou por má intenção. De considerar passos intermediários
em deduções de equações como sendo triviais, ou trivializar discussões
sobre hipóteses e fundamentos. Tentarei partir sempre de primeiros
princípios, e explicar da maneira mais simples, no mínimo passando a
receita para deduções em detalhes, indicando outras referências 
técnicas que possam complementar os textos.
